% This text is a supplement of CINTA.
% May. 2018
% Author: Libin Wang
% SCNU
% additional use of \usepackage{beamerthemesplit}
\documentclass{beamer}

%%%%%%%%%%%%%%%%%%%%%%%%%%%%%%%%%%%%%%%%%%%%%%%%%%%%%%%%%%%%%%%%%%%%%
\usepackage{xeCJK}
%\usepackage[CJKbookmarks]{hyperref}

%% Support Chinese with TeX Live fonts.

\setCJKmainfont[
  BoldFont=FandolHei-Regular.otf,
  ItalicFont=FandolKai-Regular.otf,
  BoldItalicFont=FandolKai-Regular.otf
]{FandolSong-Regular.otf}
\setCJKsansfont[AutoFakeBold, AutoFakeSlant]{FandolHei-Regular.otf}
\setCJKmonofont{FandolFang-Regular.otf}

%\setCJKmainfont[BoldFont=Microsoft YaHei,ItalicFont=YouYuan]{SimSun} %设置默认中文字体

\setmainfont{Times New Roman} % 设置默认英文衬线字体

\setCJKfamilyfont{ls}{FandolSong-Regular.otf} %定义字体
%%%%%%%%%%%%%%%%%%%%%%%%%%%%%%%%%%%%%%%%%%%%%%%%%%%%%%%%%%%%%%%%%%%%%
%\usepackage{beamerthemesplit} % new 
\usepackage{beamerthemeshadow}
\setbeamertemplate{headline}{}
\setbeamertemplate{navigation symbols}{}
\usepackage{amsmath, amsfonts, amssymb}

%% Setup for Code
\usepackage{listings}
\lstset{language = Python,
	showstringspaces = false,
	columns = fullflexible,
	numbers = left,
	numberstyle = \tiny,
	frame = single}

%%\newproof{pf}{Proof} 
%\newtheorem{thm}{Theorem}[section]
%\newtheorem{fact}[theorem]{Fact}
\newtheorem{prop}[theorem]{Proposition}
\newtheorem{lem}[theorem]{Lemma}
\newtheorem{defn}[theorem]{Definition}
\newtheorem{cor}[theorem]{Corollary}
\newtheorem{rema}[theorem]{Remark}
\newtheorem{clam}[theorem]{Claim}
\newtheorem{exm}[theorem]{Example}

%\title{A Concrete Introduction to Number Theory and Algebra -- Chapter 3-4} 
\title{Mathematics for Computer Science -- 2} 
\author{Libin Wang} 
\institute{School of Computer Science, South China Normal University}
\date{\today} 

\begin{document}

\frame{\titlepage} 

\frame{\frametitle{Table of contents}\tableofcontents} 

\section{Congruence} 
\frame{\frametitle{Definition of congruence（同余）.} 
“同余”这个概念是CSers的好朋友。
	\begin{definition}
		We say that \textit{$a$ is congruent to $b$ modulo $m$}, and we write 
		\[a \equiv b \;(\text{mod}\; m),\]
		if $m$ divides $a - b$. The number $m$ is called the \textit{modulus}（模数） of the congruence.
	\end{definition}

}

\frame{\frametitle{Notations of congruence（同余的表达）.}

\begin{exampleblock}{Some notations.} 
	$a \equiv b \pmod{m} \iff$ 
\begin{itemize}
\item $\exists k\in \mathbb{Z}$, $a = km + b$. 

\item $(a \bmod m) = (b \bmod m)$

\item $m \mid (a - b)$
\end{itemize}
\end{exampleblock}

\pause
\begin{block}{提示.}

	表达是思维的指针。
	
\end{block}
}

\frame{\frametitle{Congruence--例.} 

\begin{exm}
\[26 \equiv 8 \pmod{9} \qquad \text{and}\qquad 6 \equiv 55 \pmod{7},\]
since 
\[9 | (26 - 8)\qquad \text{and}\qquad 7 | (6 - 55),\]
or, equivalently:
\[ 8 = 26 - 2 \cdot 9 \qquad \text{and}\qquad 55 = 6 + 7 \cdot 7.\]
\end{exm}

}

\frame{\frametitle{Properties of Congruence.} 
    \begin{lem}
       	If $a_1 \equiv b_1 \pmod{m}$ and $a_2 \equiv b_2 \pmod{m}$, then
       
       \[a_1 \pm a_2 \equiv b_1 \pm b_2 \pmod{m}\]
       
       and
       
       \[a_1  a_2 \equiv b_1 b_2 \pmod{m}\]
    \end{lem}
课堂练习！
}

\frame{\frametitle{Properties of Congruence.} 
	\begin{exampleblock}{练习}
		 假设$x \equiv a \pmod{n}$，且$n = p \cdot q$，$p, q$是两个素数。请证明，因此必然有$x \equiv a \pmod{p}$和$x \equiv a \pmod{q}$。
	\end{exampleblock}
}

\frame{\frametitle{Modular Arithmetic(模算术).} 
	\begin{exampleblock}{Examples}
			Since $10000 \equiv 1 \pmod{3}$ and $998 \equiv 2 \pmod{3}$, then
			\[ 10000 \cdot 998 \equiv 2 \pmod{3}.\]
	\end{exampleblock}
}

\frame{\frametitle{Modular Arithmetic--Exercises.} 
	\begin{exampleblock}{练习}
		计算以下等式：
		\[7^{500} \bmod 8\]
		\[ 1! + 2! + 3! + \cdots + 99! + 100! \bmod 12\]
		\[2^{20} - 1 \bmod 41\]
	\end{exampleblock}
	
	\pause
	
	\begin{block}{答案}
    $1$ 、$9$、 $0$。
	\end{block}
}


\frame{\frametitle{Properties of Congruence.} 
	\begin{block}{Negative number.}
	Let $x$ and $n$ be two positive integers and $x < n$, what does $-x \bmod n$ mean?
	\end{block}\pause

    \begin{exampleblock}{Intuition.}
     Consider that the negative of $x$ is the number $x'$ such that $x + x' = 0$, that is:
     \[x + x' \equiv 0 \pmod{n}.\]
     Since $x < n$, hence $x' = n - x$.
   \end{exampleblock}
}

\frame{\frametitle{Two's Complement(二进制补码)}

\begin{block}{Two's Complement}
A signed number represented in $n$ bits. The range of the numbers is $[-2^{n-1}, 2^{n-1} - 1]$, and the rule is described as follows:
\begin{itemize}
	\item Positive integers, in the range $0$ to $2^{n-1} -1$, are stored in regular binary form. The sign bit is set to $0$.
	
	\item Negative integers $-x$, with $1 \leq x \leq 2^{n-1}$, are calculated by first constructing $x$ in binary, then inverting all the bits of $x$ and finally adding $1$. The sign bit is set to $1$.
\end{itemize}
\end{block}	
}

\frame{\frametitle{Two's Complement}
	\begin{exampleblock}{Another way to remember two's complement.}
		The negative of $x$ equals $ 2^n - x$. Since $2^n - x = \underbrace{111\cdots 11}_{n}  + 1  - x$, and $\underbrace{111\cdots 11}_{n} - x$ is the same as inverting all the bits of $x$.
	\end{exampleblock}
}

\frame{\frametitle{Cancellation Law.} 
\begin{theorem}[Cancellation Law.]
    If $\gcd(c, m) = 1$ and
    \[ac \equiv bc \pmod{m},\]
    then
    \[a \equiv b \pmod{m},\]
	where $\gcd$ shorts for the \textit{greatest common divisor}.
\end{theorem}
}

\frame{\frametitle{Cancellation Law.} 
\begin{proof} 
	By definition of  congruence, we have $m | (ac - bc)$, equivalently, $m | (a - b)c$. Since $\gcd(c,m) = 1$, it follows that $m \mid (a - b)$, so as claimed.

\end{proof}
}

\frame{\frametitle{Cancellation Law.} 
	\begin{exampleblock}{Another perspective.}
		If $\gcd(c, m) = 1$ then $\exists r,s \in \mathbb{Z}$ s.t. 
		\[r c + s m = 1\]
		both sides of the equation modulo $m$, we have:
		\[r c \equiv 1 \pmod{m}\]
		means $r$ is the multiplicative inverse(乘法逆元) of $c \bmod m$, let it be $c^{-1}$.
	\end{exampleblock}
}

\frame{\frametitle{Partially solve the congruence $a x \equiv b \pmod{m}$.}  

\begin{exm}
To solve $3 x \equiv 2 \pmod{11}$. 

Firstly, by using $e\gcd$ algorithm, to compute that $3^{-1} = 4$, because $3 \cdot 4 \equiv 1 \pmod{11}$. Multiply $4$ to the equation and obtain
\[x \equiv 8 \pmod{11}.\]
\end{exm}
}

\frame{\frametitle{Cancellation Law-- Revisited.} 
	\begin{alertblock}{问题.}
		消去律源自于同余的定义与欧几里得引理。可以执行消去操作是因为存在唯一乘法逆元，乘法逆元源自于Bezout定理，但是Bezout系数不唯一。请问，为什么乘法逆元会唯一呢？
	\end{alertblock}
	\pause 
	\begin{exampleblock}{回答.}
			需要重新审视Bezout系数的形式。
	\end{exampleblock}
}
\iffalse
\frame{\frametitle{考察Bezout系数的形式.} 
	\begin{exampleblock}{回归Bezout定理.}
		If $\gcd(a, b) = 1$ then $\exists r, s \in \mathbb{Z}$ s.t. 
		\[ar + b s = 1\]\pause
		从以上形式出发，现在想构造出新的系数应该怎么做呢？\\ \pause
		初始的想法：存在一个未知量$e$
		\[ a r + bs + e - e = 1\]		\pause
		然后分别把$e$与$ar$和$bs$“结合”起来：
		\[ ar + e + bs -e = 1\] 
	\end{exampleblock}
}

\frame{\frametitle{考察Bezout系数的形式--继续.} 
	\begin{exampleblock}{回归Bezout定理.}
		当然我们希望$e$与$ar$(或者$bs$)有共同的因子使得：
		\[ a(r + e') + b(s - e') = 1\] 
		请问，此时，我们判断$e$应该具有什么形式？\\ 
		\pause
		难道不是$e$必然同时具有$a$和$b$这两个分量吗？即$e = kab$，$k$是任意整数。把$e$带入最原始的方程得：
		\[ a r + kab + bs - kab = a(r +kb) + b(s - ka) = 1\]
		也就是说，我们构造出了一系列的Bezout系数$r + kb$和$s -k a$。
	\end{exampleblock}
}

\frame{\frametitle{考察Bezout系数的形式--继续.} 
\begin{alertblock}{问题.}
		给定$a$和$b$，$r + kb$和$s -k a$确实是Bezout系数，即满足
			\[ a(r +kb) + b(s - ka) = 1\]
		但是所有的Bezout系数都形如$r + kb$和$s -k a$吗？
	\end{alertblock}
}

\frame{\frametitle{考察Bezout系数的形式--继续.} 
	\begin{exampleblock}{思路.}
		如果同时存在两组Bezout系数满足：
		\[ a r_0 + b s_0 = 1\]
		\[ a r_1 + b s_1 = 1\]
		请问，此时，我们如何把$r_0$(或者$s_0$)表达成$r_1$与$b$（或者$s_1$与$a$）组成的形式？\\ 
		\pause
		回忆中学的“消元法”！
	\end{exampleblock}
}

\frame{\frametitle{考察Bezout系数的形式--继续.} 
	\begin{exampleblock}{思路.}
		式子一左右两边同乘$s_1$，式子二左右同乘$s_0$，所得相减：
		\[ a r_0 s_1 - a r_1 s_0 = a(r_0 s_1 - r_1 s_0) = s_1 - s_0\]
		同样的操作：
		\[b s_0 r_1 - b s_1 r_0 = b(s_0 r_1 - s_1 r_0) = r_1 - r_0\]
		\pause 
		令$k = r_0 s_1 - r_1 s_0$，可得：
		\[r_0 = r_1 + kb\]
		\[s_0 = s_1 - ka \]
	\end{exampleblock}
}
\fi
\frame{\frametitle{Properties of Congruence.} 
	\begin{lem}
		For $n \in \mathbb{N}$, congruence modulo $n$ forms an equivalence relation(等价关系) of $\mathbb{Z}$.
	\end{lem}	
	
	\begin{proof}
		It is easy to check that: 
		
		1. Reflexive(自反性). $a \equiv a \pmod{n}$
		
		2. Symmetric(对称性). If $a \equiv b \pmod{n}$ then $b \equiv a \pmod{n}$
		
		3. Transitive(传递性).	If 	$a \equiv b \pmod{n}$ and $b \equiv c \pmod{n}$ then $a \equiv c \pmod{n}$
	\end{proof}
}

\frame{\frametitle{equivalence relation.} 
	\begin{block}{问题}
		等价关系意味什么？
	\end{block}	
	\pause
	\begin{block}{条件反射式回答.}
		存在等价类，即对数据的划分(Partition)。
	\end{block}
}

\frame{\frametitle{Equivalence relation and equivalence classes.} 
\begin{defn}
When a set $\mathbb{S}$ has an equivalence relation on it, then the equivalence relation partitions(划分、分割) the set $\mathbb{S}$ into disjoint subsets, called equivalence classes（等价类）, defined by the property that two elements are in the same equivalence class if they are equivalent.

\end{defn}
}

\frame{\frametitle{Congruence classes modulo $m$.} 
The set of congruence classes modulo $m$ is denoted by $\mathbb{Z}/m\mathbb{Z}$.
There are exactly m congruence classes in $\mathbb{Z}/m\mathbb{Z}$. That is :
\[\mathbb{Z}/m\mathbb{Z} = \{[0]_m, [1]_m, \cdots, [m-1]_m\}.\]

\begin{exm}
When $m = 2$, $\mathbb{Z}/2\mathbb{Z} = \{[0]_2, [1]_2\}$. The congruence class $[1]_2$ is the set of all integers congruent to $1$ modulo $2$. Thus $[1]_2$ is the set of all odd integers. Similarly, the congruence class $[0]_2$ is the set of all even integers.
\end{exm}
}

\frame{\frametitle{Proposition of Congruence.}
\begin{prop}
\label{prop_residue}
If $[a_1]_m = [a_2]_m$ and $[b_1]_m = [b_2]_m$, then
\[[a_1 \pm b_1]_m = [a_2 \pm b_2]_m, \qquad \text{and}\qquad
[a_1 b_1]_m = [a_2 b_2]_m.\]
\end{prop}

\begin{proof}
It is easy. Transform the form $[a]_m = [b]_m$ to $a \equiv b \pmod{m}$ and use Proposition 2.2. 
\end{proof}
}

\frame{\frametitle{Notations of Congruence classes modulo $m$.} 
\begin{itemize}
\item Any element $b$ of a congruence class $[a]_m$ is called a \emph{representative} (代表元) of that class. 

\item The set of all the least nonnegative representative of $\mathbb{Z}/m\mathbb{Z}$ is the set of integers $\{0, 1, 2,\cdots, m - 1 \}$, that is called the \textit{least residue system} modulo $m$ (模$m$的最小剩余系). 

\item Any set of $m$ integers, no two of which are congruent modulo $m$, is called a \textit{complete residue system} modulo $m$ (模$m$的完全剩余系).
\end{itemize}

}

\frame{\frametitle{Notations of Congruence classes modulo $m$.} 

\begin{exm}
Let $m = 7$, the least residue system modulo $m$ is the set $\{0, 1, 2, 3, 4, 5, 6 \}$, 
and a complete residue system modulo $m$ may be the set $\{14 , 8, 23, 46, 61, 13\}$.
\end{exm}
}

\section{Modular Exponentiation}
\label{sec:mod-pow}
\frame{\frametitle{Mod Exponentiation.}
In this section, we focus on modular exponentiation which is an important arithmetic primitive. Its task is that given integers $x$, $y$ and $m$ to compute 
\[ x^ y \bmod m.\] 
}
 
\frame{\frametitle{Mod Exponentiation.}
\begin{exm}
\label{exm:pow}
To compute $2^{16} \bmod 11$.  We compute:
\[2^2 \bmod 11 = 4\] 
\[2^4 \bmod 11 = 4 \cdot 4 \bmod 11 = 5\]
\[2^8 \bmod 11 = 5 \cdot 5 \bmod 11 = 3 \]
\[2^{16} \bmod 11 = 3 \cdot 3 \bmod 11 = 9 \]
\end{exm}
}

\frame{\frametitle{Mod Exponentiation.}
\begin{exm}
\label{exm:pow-22}
To compute $2^{22} \bmod 11$.  We compute:
\[2^2 \bmod 11 = 4\] 
\[2^4 \bmod 11 = 4 \cdot 4 \bmod 11 = 5\]
\[2^8 \bmod 11 = 5 \cdot 5 \bmod 11 = 3 \]
\[2^{16} \bmod 11 = 3 \cdot 3 \bmod 11 = 9 \]

then \[(2^{22} = 2^2 \cdot 2^4 \cdot 2^{16}) \equiv 4 \cdot 5 \cdot 9 \equiv 4 \pmod{11}\]
\end{exm}
}

\frame{\frametitle{Mod Exponentiation.}
The process can be expressed as a recursive form, by that, we sharply improve the efficiency from performing $O(y)$ multiplications to $O(log (y))$.
\begin{equation}  
\label{eq:rec_pow}
x ^ y = 
\left\{  
\begin{array}{lr}  
(x ^{ \lfloor y/2 \rfloor })^2 &  \mbox{if $y$ is even;}\\  

x \cdot (x ^{ \lfloor y/2 \rfloor })^2 & \mbox{if $y$ is odd.} \\  

\end{array}  
\right.  
\end{equation}
}


\frame{\frametitle{Mod Exponentiation(Recursive Version).}
\begin{center}
\label{alg:rec_pow}
	\lstinputlisting[caption = Recursive Modular Exponentiation]{../codes/chap2-3.py}
\end{center}
}

\frame{\frametitle{Mod Exponentiation: from recursive to iterative.}
We describe how to transform the recursive algorithm to an iterative algorithm as follows. 
Firstly, we treat integer $y$ as a polynomial (or a binary string):
\[ y = y_{n-1}2^{n-1} + y_{n-2}2^{n-2} \cdots + y_1 2 + y_0 ,\] where $y_i \in \{0, 1\}$. \pause

Secondly, transform $x^y$ as:
\[x^y = \prod_{i=0}^{n-1} x^{y_i 2^{i}}\] \pause

Finally, start with x and repeatedly square modulo $m$,  multiply the terms with $y_i = 1$ and get the result.
}

\frame{\frametitle{Mod Exponentiation: Example.}

\begin{exampleblock}{A comcrete example.}
	Let $x=7$, $y= 10$, $m = 11$,
	to compute $x^y \bmod m$. The bianry string of $y$ is $1010$, thus we compute:
	\begin{align*}
		y_0 &= 0, & x^{2^0} &\equiv 7 \pmod{m}\\
		y_1 &= 1, & x^{2^1} &\equiv 5 \pmod{m}\\
		y_2 &= 0, & x^{2^2} &\equiv 3 \pmod{m}\\
		y_3 &= 1, & x^{2^3} &\equiv 9 \pmod{m}
	\end{align*}
Then, multiply the terms with $y_i = 1$, we have $x^y = (5 \cdot 9) \bmod 11 = 1$
\end{exampleblock}
}

\frame{\frametitle{Mod Exponentiation: Exercise.}

\begin{block}{练习.}
	Let $x=5$, $y= 13$, $m = 13$,
	to compute $x^y \bmod m$. 
\end{block}
\pause

\begin{block}{答案.}
5. 
\end{block}
}

\frame{\frametitle{Mod Exponentiation (Iterative version).}

Because the white board is too narrow to show the code, so it is your home work.
\iffalse
\begin{center}
	\lstinputlisting[caption = Iterative Modular Exponentiation]{../codes/chap2-4.py}
\end{center}
\fi
}

\section{Fermat's Little Theorem} 

\frame{\frametitle{Some topics using modular arithmetic.}

Our following job is to play with number using modular arithmetic, and find some patterns or rules.

}

\frame{\frametitle{Find the patterns.}
Let $p = 7$, and for very $1\leq a < p$, compute $a^i\; \text{mod}\; p$, where $1\leq i < p$. We have:
\begin{center}
	\begin{tabular}{ | c | c | c | c | c | c | }
		\hline
		$a$  & $a^2$ & $a^3$ & $a^4$ & $a^5$ & $a^6$ \\ \hline
		1 & 1 & 1 & 1 & 1 & 1 \\ \hline
		2 & 4 & 1 & 2 & 4 & 1 \\ \hline
		3 & 2 & 6 & 4 & 5 & 1 \\ \hline
		4 & 2 & 1 & 4 & 2 & 1 \\ \hline
		5 & 4 & 6 & 2 & 3 & 1 \\ \hline
		6 & 1 & 6 & 1 & 6 & 1 \\
		\hline
	\end{tabular}
\end{center}\pause 

May you find some patterns?
}

\frame{\frametitle{More data to find the patterns.}
	Let $p = 11$, and for very $1\leq a < p$, compute $a^i\; \text{mod}\; p$, where $1\leq i < p$. We have:
	\begin{center}
		\begin{tabular}{ | c | c | c | c | c | c | c | c | c | c |}
			\hline
			$a$  & $a^2$ & $a^3$ & $a^4$ & $a^5$ & $a^6$ & $a^7$ & $a^8$ & $a^9$ & $a^{10}$ \\ \hline
			1 & 1 & 1 & 1 & 1 & 1 & 1 & 1 & 1 & 1\\ \hline
			2 & 4 & 8 & 5 & 10 & 9 & 7 & 3 & 6 & 1 \\ \hline
			3 & 9 & 5 & 4 & 1 & 3 & 9 & 5 & 4 & 1 \\ \hline
			4 & 5 & 9 & 3 & 1 & 4 & 5 & 9 & 3 & 1 \\ \hline
			5 & 3 & 4 & 9 & 1 & 5 & 3 & 4 & 9 & 1 \\ \hline
			6 & 3 & 7 & 9 & 10 & 5 & 8 & 4 & 2 & 1 \\ \hline
			7 & 5 & 2 & 3 & 10 & 4 & 6 & 9 & 8 & 1 \\ \hline
			8 & 9 & 6 & 4 & 10 & 3 & 2 & 5 & 7 & 1 \\ \hline
			9 & 4 & 3 & 5 & 1 & 9 & 4 & 3 & 5 & 1 \\ \hline
			10 & 1 & 10 & 1 & 10 & 1 & 10 & 1 & 10 & 1  \\ 
			\hline
		\end{tabular}
	\end{center}
	
}

\frame{\frametitle{Conjecture.}
\begin{block}{猜想}
For $p$ is a prime, and $\gcd(a,p) = 1$,  
\[a^{p-1} \equiv 1 \pmod{p}\]
\end{block}
}


\frame{\frametitle{Another computation.}
$\forall 1 < a < p$, compute $a \cdot i \; \text{mod} \; p$, for $1 \leq i < p$. For example, let $a = 2$, $p = 7$, we have: 
\begin{center}
	\begin{tabular}{ | c | c | c | c | c | c | c | }
		\hline
		 $a \cdot i$  & 1 & 2 & 3 & 4 & 5 & 6 \\ \hline
		 $a=1$ & 1 & 2 & 3 & 4 & 5 & 6 \\ \hline 
		 $a=2$ & 2 & 4 & 6 & 1 & 3 & 5 \\ \hline 
		 &  &  &  &  & &  \\ \hline 
		 &  &  &  &  & & \\ \hline
         &  &  &  &  & & \\ \hline
		 &  &  &  &  & & \\ 
		\hline
	\end{tabular}
\end{center}
}

\frame{\frametitle{Another computation.}
 Continue the computation...
	\begin{center}
		\begin{tabular}{ | c | c | c | c | c | c | c | }
			\hline
			$a \cdot i$  & 1 & 2 & 3 & 4 & 5 & 6 \\ \hline
			$a=1$ & 1 & 2 & 3 & 4 & 5 & 6 \\ \hline 
			$a=2$ & 2 & 4 & 6 & 1 & 3 & 5 \\ \hline 
			$a=3$ & 3 & 6 & 2 & 5 & 1 & 4 \\ \hline 
			&  &  &  &  & & \\ \hline
			&  &  &  &  & & \\ \hline
			&  &  &  &  & & \\ 
			\hline
		\end{tabular}
	\end{center}\pause
Have a decision?
}

\frame{\frametitle{Another computation.}
	Finally!
	\begin{center}
		\begin{tabular}{ | c | c | c | c | c | c | c | }
			\hline
			$a \cdot i$  & 1 & 2 & 3 & 4 & 5 & 6 \\ \hline
			$a=1$ & 1 & 2 & 3 & 4 & 5 & 6 \\ \hline 
			$a=2$ & 2 & 4 & 6 & 1 & 3 & 5 \\ \hline 
			$a=3$ & 3 & 6 & 2 & 5 & 1 & 4 \\ \hline 
			$a=4$& 4 & 1 & 5 & 2 & 6 & 3 \\ \hline
			$a=5$& 5 & 3 & 1 & 6 & 4 & 2 \\ \hline
			$a=6$& 6 & 5 & 4 & 3 & 2 & 1 \\ 
			\hline
		\end{tabular}
	\end{center}
}

\frame{\frametitle{Conjecture.}
The trick is, $p$ is a prime number!
We conjecture:
if $p$ is a prime, $\forall a$ which is not divided by $p$, $a \cdot i \bmod p$, for $1 \leq i < p$, is a permutation of numbers from $1$ to $p-1$.\pause
That is,  the numbers
\[a, 2a, 3a, \cdots, (p-1)a \pmod{p}\]
are the same as the numbers:
\[1, 2, 3, \cdots, p-1\]
although they may be in a different order. \pause
\[S  = \{ a \cdot i \bmod p, 1 \leq i < p \}\]
is also called a complete system of residues modulo $p$.
}

\frame{\frametitle{Proof by contradiction.}
	Of course, we need a proof!\pause
	\begin{proof} {Proof by contradiction (informal and incomplete).} 
		If we are wrong, then there exist $i$ and $j$ such that,
		\[a \cdot i \equiv a \cdot j \pmod{p}\]
		where $i \neq j$.
		However, then we can cancel the $a$ from the equation! (Cancellation Law.)
	\end{proof}
	
}

\frame{\frametitle{Do a simple job!}
	Multiply all $1 \leq i < p$, and all $a \cdot i \bmod p$ , we have:
	
	\[\prod_{i=1}^{p-1} i = \prod_{i=1}^{p-1} (a \cdot i \bmod p)\]
	
	\pause
	 Convince yourself:
	\[\prod_{i=1}^{p-1} i \equiv \prod_{i=1}^{p-1} a \cdot i \equiv a^{p-1}\prod_{i=1}^{p-1} i \pmod{p}\]  \pause
	Cancel the big number, we have:
	\[a^{p-1} \equiv 1 \pmod{p}\]
}

\frame{\frametitle{Fermat's little theorem.}
	\begin{theorem}[Fermat's little theorem.]
	Let $p$ be a prime number, and let $a$ be any number with $a \not \equiv 0 \pmod{p}$. Then
	\[a^{p-1} \equiv 1 \pmod{p}.\]
	\end{theorem}
}

\frame{\frametitle{An Exercise using Fermat's little theorem.}
	\begin{exampleblock}{Exercise.}
	Let $p = 17$ be a prime number, and let $a = 2$ , what is $a^{2024} \bmod p$? \pause
	
	\[2^{2024} \equiv 2^{126 \cdot 16 + 8 } \equiv 2^8 \equiv 1  \pmod{17}.\]
	\end{exampleblock}
}

\frame{\frametitle{Exercises using Fermat's little theorem.}
	\begin{exampleblock}{Exercises.}
	\begin{itemize}
	\item $a \equiv 9^{794} \pmod{73}$, find $a$ and $0 \le a \le 73$.
	
	\item Solve $x^{87} \equiv 6 \pmod{29}$.
	\end{itemize}
	\end{exampleblock}
}

\section{Euler's Theorem} 

\frame{\frametitle{More computation.}
	If the modulus is a composite number, then our trick will fail! For example, let $n = 6$, 
		\begin{center}
		\begin{tabular}{ | c | c | c | c | c | c |}
			\hline
			$a \cdot i$  & 1 & 2 & 3 & 4 & 5 \\ \hline
			$a=1$ & 1 & 2 & 3 & 4 & 5 \\ \hline 
			$a=2$ & 2 & 4 & 0 & 2 & 4 \\ \hline 
			$a=3$ & 3 & 0 & 3 & 0 & 3 \\ \hline 
			$a=4$ & 4  & 2 & 0 & 4 & 2 \\ \hline
			$a=5$ &  5 & 4  & 3  & 2  & 1  \\ 
			\hline
		\end{tabular}
	\end{center}

}

\frame{\frametitle{One more computation.}
Let $n = 9$,
	\begin{center}
	\begin{tabular}{ | c | c | c | c | c | c | c | c | c | }
		\hline
		$a \cdot i$  & 1 & 2 & 3 & 4 & 5 & 6 & 7 & 8 \\ \hline
		$a=1$ & 1 & 2 & 3 & 4 & 5 & 6 & 7 & 8 \\ \hline 
		$a=2$ & 2 & 4 & 6 & 8 & 1  & 3 & 5 & 7 \\ \hline 
		$a=3$ & 3 & 6 & 0 & 3 & 6 & 0 & 3 & 6 \\ \hline 
		$a=4$ & 4 & 8 & 3 & 7 & 2 & 6 & 1 & 5 \\ \hline
		$a=5$ & 5 & 1 & 6  & 2  & 7 & 3 & 8 & 4 \\ \hline
		$a=6$ & 6 & 3 & 0 & 6 & 3 & 0 & 6 & 3 \\ \hline 
		$a=7$ & 7 & 5 & 3 & 1 & 8 & 6 & 4 & 2 \\ \hline
		$a=8$ & 8 & 7  & 6  & 5  & 4 & 3 & 2 & 1 \\ 
		\hline
	\end{tabular}
\end{center}
}

\frame{\frametitle{Observation.}
\begin{block}{观察.}	
	If $n$ is a composite number, the numbers
	\[a, 2a, 3a, \cdots, (n-1)a \pmod {n}\]
	may \textbf{NOT} the same as the numbers:
	\[1, 2, 3, \cdots, n-1\]
\end{block}
}

\frame{\frametitle{Check the observation.}
	Let $n = 9$，for all $i \in [1, n-1]$，and $\gcd(i, n) =1$: 
	\begin{center}
		\begin{tabular}{ | c | c | c | c | c | c | c | }
			\hline
			$a \cdot i$  & 1 & 2 & 4 & 5 & 7 & 8 \\ \hline
			$a=1$ & 1 & 2 & 4 & 5 & 7 & 8 \\ \hline 
			$a=2$ & 2 & 4 & 8 & 1 & 5 & 7 \\ \hline 
			$a=4$ & 4 & 8 & 7 & 2 & 1 & 5 \\ \hline
			$a=5$ & 5 & 1 & 2 & 7 & 8 & 4 \\ \hline
			$a=7$ & 7 & 5 & 1 & 8 & 4 & 2 \\ \hline
			$a=8$ & 8 & 7 & 5 & 4 & 2 & 1 \\ 
			\hline
		\end{tabular}
	\end{center}	
}

\frame{\frametitle{Conjecture.}
	We conjecture:
	Let $n$ be a composite number, denotes
	\[S = \{ b: 1\leq b < n \; \text{and}\; \gcd(b,n)=1 \}\]
	Then $\forall a$ with $\gcd(a,n)=1$, denotes
	\[S' = a \cdot S  \pmod{n}\] 
	we have:
	\[S = S'\]
}

\frame{\frametitle{Notation.}
	\begin{block}{Euler's phi function.}
	Define:
	\[\varphi(n) = | \{ b: 1\leq b < n \; \text{and}\; \gcd(b,n)=1 \}|\]
	The function $\varphi$ is called \textit{Euler's phi function}.
	\end{block}
}

\frame{\frametitle{Notation.}
	Then:
	\[
	S =  \{ b_1, b_2, \cdots ,  b_{\varphi(n)}:
	1\leq b_i < n \; \text{and}\; \gcd(b_i,n)=1 \}
	\]
	\[
	\begin{split}
	S' =  \{& a \cdot b_1, a \cdot b_2, \cdots , a \cdot b_{\varphi(n)}\pmod{n}:\\
	& b_i \in S \; \text{and}\; \gcd(a,n)=1 \}.
	\end{split}
	\]
}

\frame{\frametitle{Proof.}
	To prove
	\[S = S'\]
	
	Check that if there exist: 
	\[a \cdot b_i \equiv a \cdot b_j \pmod  n\]
	where $b_i \neq b_j$. Then by Cancellation Law,
    \[b_i \equiv b_j  \pmod  n.\] 
    Contradiction!
}

\frame{\frametitle{Do a similar simple job!}
	Multiply all the numbers in $S$ and $S'$ , we have:
	
	\[\prod_{i = 1}^{\varphi(n)} b_i = \prod_{i=1}^{\varphi(n)} (a \cdot b_i \bmod n)\]

	\[\prod_{i = 1}^{\varphi(n)} b_i \equiv \prod_{i=1}^{\varphi(n)} a \cdot b_i \pmod  n\]
	
	\[\prod_{i = 1}^{\varphi(n)} b_i \equiv a^{\varphi(n)} \prod_{i=1}^{\varphi(n)} b_i \pmod n\]
	
	Cancel the big number, we have:
	\[a^{\varphi(n)} \equiv 1 \pmod n\]
	
}

\frame{\frametitle{Euler's Theorem.}
	\begin{block}{(Euler's Theorem.)}
	Let $n$ be a positive composite number, $a$ be a positive integer with $\gcd(a, n) = 1$, then:
	\[a^{\varphi(n)} \equiv 1 \pmod{n}.\]
	
\end{block}
}


\frame{\frametitle{Example.}

    \begin{exampleblock}{例.}
    	设$n =  100$, 则$\varphi(n) = 40$, 所以$7^{40} \equiv 1 \pmod{n}$。同样，$9^{40} \equiv 1 \pmod{n}$。  
    \end{exampleblock}
    \pause
    \begin{alertblock}{问题.}
    为什么$\varphi(n) = 40$？或者，给定$n$，如何求$\varphi(n)$？
    \end{alertblock}

}

\frame{\frametitle{More about Euler's Phi Function.}
	\begin{defn}[Euler's Phi Function.]
		Define:
		\[\varphi(n) = | \{ b: 1\leq b < n \; \text{and}\; \gcd(b,n)=1 \}|\]
		The function $\varphi$ is called \textit{Euler's phi function}.
	\end{defn}

    \begin{exampleblock}{Observations}
    	$\varphi(p) = p - 1$, where $p$ is a prime.
    	
    	$\varphi(p^k) = p^k - p^{k-1}$, where $p$ is a prime. 
    \end{exampleblock}

}

\frame{\frametitle{More about Euler's Phi Function..}
	
	\begin{exampleblock}{Question.}
		How to compute $\varphi(m)$ where $m = p^i q^j$ with $p$ and $q$ are prime. 
	\end{exampleblock}\pause

\begin{exampleblock}{A related question.}
	How to compute $\varphi(mn)$ where $\gcd(m,n) = 1$. 
\end{exampleblock}
	
}

\frame{\frametitle{More about Euler's Phi Function.}

	\begin{exampleblock}{Compute $\varphi(mn)$}
Display the positive integers not exceeding $mn$ in the following way.

		\begin{tabular}{c c c c c}
		$1$ & $m+1$  & $2m+1$ & $\cdots$ & $(n-1)m+1$ \\ 
		$2$ & $m+2$ & $2m+2$ & $\cdots$ & $(n-1)m+2$ \\ 
		$3$ & $m+3$ & $2m+3$ & $\cdots$& $(n-1)m+3$ \\ 
		$\cdots$ & $\cdots$ & $\cdots$ & $\cdots$ & $\cdots$ \\ 
		$r$ & $m+r$ & $2m+r$ & $\cdots$  &$(n-1)m+r$  \\  
		$\cdots$&  $\cdots$ & $\cdots$ & $\cdots$ & $\cdots$ \\ 
		$m$ & $2m$ & $3m$  & $\cdots$ & $mn$ \\ 
		\end{tabular} 
	
	\end{exampleblock}
}

\frame{\frametitle{More about Euler's Phi Function.}
\begin{block}{Basic idea.}
Find all the elements which are relatively prime to both $n$ and $m$, then it is relatively prime to $mn$. Formally:
\[ \forall a, \gcd(a, n) = 1 \; \text{and} \; \gcd(a, m) = 1 \implies \gcd(a, mn) = 1. \]
\end{block}
}

\frame{\frametitle{More about Euler's Phi Function.}
\begin{block}{Counting}
\begin{itemize}
	\item How many rows satisfy $\gcd(r,m) =1$? Ans : $\varphi(m)$. Note that, if  $\gcd(r,m) =1$ then $\gcd(km + r, m) =1$, for $k \in [0..n-1]$.
	
	\item At $r$th row, how many integers have $\gcd(km + r, n)=1$, for $k \in [0..n-1]$? Ans: $\varphi(n).$ Note that, we have $\gcd(n,m)=1$.
	
	\item Hence, there are $\varphi(m)$ rows, each containing $\varphi(n)$ integers relatively prime to $mn$.
\end{itemize}
\end{block}
}

\frame{\frametitle{More about Euler's Phi Function.}
\begin{rema}[Why the conclusion holds?]
   If $\gcd(a, m)=1$ and $\gcd(a, n)=1$ then $\gcd(a, mn) = 1$.
\end{rema}
}

\frame{\frametitle{More about Euler's Phi Function.}
\begin{rema}[Why the second item holds?]
   The elements in $r$th row are: $r$, $m + r$, $\cdots$, $(n-1)m + r$ with $\gcd(m, n) = 1$. $\forall k_i \neq k_j$, $k_i m + r \not \equiv k_j m + r \pmod n$. Otherwise, $k_i = k_j$ by our Golden Law(Cancellation Law), contradiction! It means the $n$ elements in $r$th row form "a complete system of residues modulo $n$", that is $\{ r, m + r, \cdots, (n-1)m + r\} \pmod n = \{0, 1, 2, \cdots, n-1\}$. Hence, exactly $\varphi(n)$ of these integers are relatively prime to $n$.
\end{rema}
}

\frame{\frametitle{More about Euler's Phi Function.}
	\begin{theorem}
		Let $m$ and $n$ be relatively prime positive integers. Then $\varphi(mn) = \varphi(m) \varphi(n)$.
\end{theorem}
}

\frame{\frametitle{More about Euler's Phi Function.}
	\begin{block}{Some easy generalizations.}
		How to relate $\varphi(mn) = \varphi(m) \varphi(n)$ where $\gcd(m,n) =1$ with $\varphi(m)$ where $m = p^i q^j$ ? \pause
		
		Ans: $\gcd(p^i, q^j) = 1$ when $p$ and $q$ are relatively prime. Hence $\varphi(p^i q^j) = \varphi(p^i) \varphi(q^j).$\pause

	   How to generalize the result to $\varphi(m)$ where $m = p_1^{a_1} p_2^{a_2}\cdots p_k^{a_k}$ ? \pause
	
    	Ans: Induction! Left as an exercise.
\end{block}
}

\frame{\frametitle{Exercise.}
	
	\begin{block}{习题.}
	
	求$a$，使得$a$满足以下条件：
	\begin{itemize}
	\item $a \equiv 7^{2003} \pmod{3750}$
	
	\item $1\le a \le 5000$
	
	\item $a$不被$7$整除.
	
	\end{itemize}
	
    \end{block}
}

\frame{\frametitle{Exercise.}

    \begin{block}{习题.}
	使用费尔马小定理求解同余方程$x^{51} \equiv 2 \pmod{17}$和$x^{50} \equiv 2 \pmod{17}$。%ans: 8、6和11
	
    \end{block}
	
	\begin{block}{习题.}
	
	%%% 113 \cdot 17 = 1 mod 192, 所以2^17 mod 221 =  19即为所求.  
	请求解$x^{113} \equiv 2 \pmod{221}$。[提示: $113$与$\varphi(221)$互素。]
	
    \end{block}
}

\frame{\frametitle{Exercise.}
	Prove the following theorem.
	\begin{theorem}[Euler's Phi Function.]
	
		Let $m = p_1^{a_1} p_2^{a_2}\cdots p_k^{a_k}$,  for all $p_i$ ($(1 \le i \le k)$) is a prime, then $\varphi(m) = m(1 - 1/p_1)(1 - 1/p_2) \cdots (1 - 1/p_k)$.
	
    \end{theorem}
}

%\iffalse

\section{RSA}

\frame{\frametitle{Public Key Cryptography.}
\begin{itemize}
	
	\item $pk, sk \leftarrow KG(\lambda)$, $sk$ is local and secret, while $pk$ is public.
	
	\item Given plaintext $m$, compute ciphertext $ c \leftarrow Enc(m, pk)$ 
	
	\item Given ciphertext $c$, compute plaintext $m \leftarrow Dec(c, sk)$
	\end{itemize}
}

\frame{\frametitle{RSA}
	\begin{itemize}
		
		\item $pk, sk \leftarrow KG(\lambda)$:
		\begin{itemize}
		\item Choose two large primes $p$ and $q$, $n = p\cdot q$, let $\varphi(n) = (p-1)(q-1).$
		
		\item Choose a integer $e$, with $\gcd(e, \varphi(n))=1$, compute $d$, s.t. $ed \equiv 1 \pmod{\varphi(n)}.$
		
		\item Let $pk$ be $e$ and $n$; $sk$ be $d$ and $n$.
		
		\end{itemize}
		\item $Enc(m, pk)$: let $c \equiv m^e \pmod{n}$
		
		\item $Dec(c, sk)$: let $m \equiv c^d \pmod{n}$
	\end{itemize}
}

\frame{\frametitle{Why RSA is correct?}
	Because of Euler's Theorem.
	\begin{itemize}	
		\item $KG(\lambda)$:
		\begin{itemize}
			
			\item $e$ and $d$ satisfy  $ed \equiv 1 \pmod{\varphi(n)}.$
			
			\item means: $\exists k$, s.t. $ed = k\varphi(n) +1 $
			
		\end{itemize}
		\item $Enc(m, pk)$: $c \equiv m^e \pmod{n}$
		
		\item $Dec(c, sk)$: $c^d \equiv m^{ed} \equiv m^{k\varphi(n)+1} \equiv m \pmod{n}$
	\end{itemize}
}

\frame{\frametitle{RSA example.}
	
	\begin{itemize}
		
		\item $KG(\lambda)$:
		\begin{itemize}
			\item Choose two primes $p=7$ and $q=11$, $n = p \cdot q= 77$, let $\varphi(n) = (p-1)(q-1)=60.$
			
			\item Choose a integer $e =7$, with $\gcd(e, \varphi(n))=1$, compute $d$, s.t. $ed \equiv 1 \pmod{\varphi(n)}$, $d = 43$.
			
			\item Let $pk$ be $e=7$ and $n=77$; $sk$ be $d=43$ and $n=77$.
			
		\end{itemize}
		\item $Enc(m,pk)$: Given $m=3$, then $c = 3^7 \; \bmod 77 = 31$.
		
		\item $Dec(c, sk)$: Let $m = 31^{43} \bmod 77 = 3$.
	\end{itemize}
}
%\fi
\section{Summary}
\frame{\frametitle{What is a Concrete Introduction?}
	\begin{itemize}
		\item Play with numbers. 
		\item Find the patterns, find the fun.
		\item Programming is a good way to play.
	\end{itemize}
}

\frame{\frametitle{What have been covered?}
	\begin{itemize}
		\item Congruence.
		\item Fermat's little theorem.
		\item Euler's theorem.
		%\item Application: a PKC algorithm RSA.
	\end{itemize}
}

\frame{\frametitle{What have been omitted?}
	\begin{itemize}
		\item Fast multiplication.
		\item Powers: how to do fast power.
	%	\item Many about RSA: security and efficiency.
	\end{itemize}
}

\frame{\frametitle{What is the next step?}
	\begin{itemize}
		\item From a new perspective to view Fermat's and Euler's theorems.
		
		\item From arithmetic go to algebra.
	\end{itemize}
}

\end{document}
